% ═══════════════════════════════════════════════════════════════
% SPANISCH LERNBLATT: Verbos regulares -ER e -IR
% ═══════════════════════════════════════════════════════════════
% Thema: Regelmäßige Verben der 2. und 3. Konjugation
% Klasse 7 | Schulkontext | 90 Minuten
% ═══════════════════════════════════════════════════════════════

\documentclass[10pt, a4paper]{article}

\usepackage[utf8]{inputenc}
\usepackage[T1]{fontenc}
\usepackage[ngerman]{babel}

% --- SCHRIFTEN ---
\usepackage{palatino}
\usepackage{helvet}
\newcommand{\sanstext}[1]{{\fontfamily{phv}\selectfont #1}}

\usepackage{microtype}
\usepackage[margin=1.4cm, top=1.6cm, bottom=1.3cm]{geometry}
\usepackage{enumitem}
\usepackage{tabularx}
\usepackage{booktabs}
\usepackage{array}
\usepackage{multicol}
\usepackage{tikz}
\usetikzlibrary{calc}

\usepackage{fancyhdr}
\usepackage{lastpage}
\usepackage{needspace}

\usepackage{xcolor}
\usepackage{amssymb}
\usepackage{tcolorbox}
\tcbuselibrary{skins}

% ═══════════════════════════════════════════════════════════════
% FARBEN
% ═══════════════════════════════════════════════════════════════

\definecolor{kopfzeile}{HTML}{1A1A1A}
\definecolor{akzent}{HTML}{C41E3A}          % Spanisch-Karminrot
\definecolor{hellgrau}{HTML}{F2F2F2}
\definecolor{mittelgrau}{HTML}{777777}
\definecolor{dunkelgrau}{HTML}{444444}
\definecolor{endungER}{HTML}{1565C0}        % Blau für -ER
\definecolor{endungIR}{HTML}{6A1B9A}        % Violett für -IR
\definecolor{endungAR}{HTML}{2E7D32}        % Grün für -AR (Wiederholung)

% ═══════════════════════════════════════════════════════════════
% VARIABLEN
% ═══════════════════════════════════════════════════════════════

\newcommand{\thema}{Verbos regulares: -ER e -IR}
\newcommand{\klasse}{7}

% ═══════════════════════════════════════════════════════════════
% KOPF- UND FUSSZEILE
% ═══════════════════════════════════════════════════════════════

\pagestyle{fancy}
\fancyhf{}
\renewcommand{\headrulewidth}{0pt}
\renewcommand{\footrulewidth}{0pt}
\cfoot{\sanstext{\footnotesize\textcolor{mittelgrau}{\thepage\,/\,\pageref{LastPage}}}}

% ═══════════════════════════════════════════════════════════════
% BOXEN
% ═══════════════════════════════════════════════════════════════

% Grammatik-Box (dunkelgrauer Strich oben)
\newtcolorbox{grammatikbox}[1][]{
    colback=hellgrau,
    colframe=hellgrau,
    boxrule=0pt,
    arc=0pt,
    left=8pt, right=8pt, top=8pt, bottom=6pt,
    borderline north={1.5pt}{0pt}{dunkelgrau},
    title={\sanstext{\small\bfseries #1}},
    coltitle=dunkelgrau,
    attach boxed title to top left={yshift=-2pt, xshift=8pt},
    boxed title style={colback=hellgrau, colframe=hellgrau, boxrule=0pt, arc=0pt}
}

% Merk-Box (roter Strich oben)
\newtcolorbox{merkbox}[1][]{
    colback=hellgrau,
    colframe=hellgrau,
    boxrule=0pt,
    arc=0pt,
    left=8pt, right=8pt, top=8pt, bottom=6pt,
    borderline north={1.5pt}{0pt}{akzent},
    title={\sanstext{\small\bfseries\textcolor{akzent}{#1}}},
    coltitle=akzent,
    attach boxed title to top left={yshift=-2pt, xshift=8pt},
    boxed title style={colback=hellgrau, colframe=hellgrau, boxrule=0pt, arc=0pt}
}

% Vokabel-Box (dünner Rahmen)
\newtcolorbox{vokabelbox}[1][]{
    colback=white,
    colframe=mittelgrau,
    boxrule=0.5pt,
    arc=0pt,
    left=8pt, right=8pt, top=8pt, bottom=6pt,
    title={\sanstext{\small\bfseries #1}},
    coltitle=dunkelgrau,
    attach boxed title to top left={yshift=-2pt, xshift=8pt},
    boxed title style={colback=white, colframe=white, boxrule=0pt, arc=0pt}
}

% AB-Verweis
\newcommand{\abmark}[1]{%
    \tikz[baseline=(X.base)]\node[fill=akzent, text=white,
    font=\sanstext\tiny\bfseries, inner sep=2pt, rounded corners=1pt](X){$\rightarrow$ AB #1};%
}

% ═══════════════════════════════════════════════════════════════
% BEFEHLE
% ═══════════════════════════════════════════════════════════════

% Spanisch: Palatino fett
\newcommand{\es}[1]{\textbf{#1}}

% Deutsch: klein, kursiv, grau
\newcommand{\de}[1]{{\small\textit{\textcolor{mittelgrau}{#1}}}}

% Grammatikbegriff: Kapitälchen
\newcommand{\gram}[1]{\textsc{\small #1}}

% Endungen farbig
\newcommand{\ender}[1]{\textcolor{endungER}{\textbf{#1}}}
\newcommand{\endir}[1]{\textcolor{endungIR}{\textbf{#1}}}
\newcommand{\endar}[1]{\textcolor{endungAR}{\textbf{#1}}}

\setlength{\parindent}{0pt}
\setlength{\parskip}{4pt}
\setlength{\columnsep}{20pt}

% ═══════════════════════════════════════════════════════════════
% DOKUMENT
% ═══════════════════════════════════════════════════════════════

\begin{document}

% ═══════════════════════════════════════════════════════════════
% HEADER
% ═══════════════════════════════════════════════════════════════

\noindent
\begin{tikzpicture}[remember picture, overlay]
    \fill[kopfzeile] (current page.north west) rectangle +(\paperwidth, -1cm);
    \node[anchor=west, text=white, font=\sanstext\large\bfseries]
        at ([xshift=1.4cm, yshift=-0.5cm]current page.north west) {\thema};
    \node[anchor=east, text=white, font=\sanstext\small]
        at ([xshift=-1.4cm, yshift=-0.5cm]current page.north east)
        {Spanisch\,·\,Klasse \klasse\,·\,Lernblatt};
\end{tikzpicture}

\vspace{0.6cm}

% ═══════════════════════════════════════════════════════════════
% INHALT
% ═══════════════════════════════════════════════════════════════

\begin{multicols}{2}

% ═══════════════════════════════════════════════════════════════
% ABSCHNITT 1: Einführung
% ═══════════════════════════════════════════════════════════════

\begin{grammatikbox}[Die drei Konjugationen]

\vspace{2pt}

Im Spanischen gibt es \textbf{drei Gruppen} von regelmäßigen Verben:

\vspace{4pt}

{\small
\renewcommand{\arraystretch}{1.3}
\begin{tabular}{@{}lll@{}}
\textbf{1. Konjugation} & \endar{-AR} & \es{hablar, escuchar, mirar} \\
\textbf{2. Konjugation} & \ender{-ER} & \es{leer, comprender} \\
\textbf{3. Konjugation} & \endir{-IR} & \es{abrir, escribir} \\
\end{tabular}}

\vspace{6pt}

{\footnotesize\textbf{Tipp:} Die Endungen von \ender{-ER} und \endir{-IR} sind \textit{fast gleich}!}

\end{grammatikbox}

\vspace{4pt}

% ═══════════════════════════════════════════════════════════════
% ABSCHNITT 2: Konjugation -ER
% ═══════════════════════════════════════════════════════════════

\begin{merkbox}[Verbos en -ER: leer, comprender]

\vspace{2pt}

{\small
\renewcommand{\arraystretch}{1.4}
\begin{tabular}{@{}l|l|l@{}}
& \textbf{leer} \de{lesen} & \textbf{comprender} \de{verstehen} \\
\hline
yo & le\ender{o} & comprend\ender{o} \\
tú & le\ender{es} & comprend\ender{es} \\
él/ella & le\ender{e} & comprend\ender{e} \\
nosotros & le\ender{emos} & comprend\ender{emos} \\
vosotros & le\ender{éis} & comprend\ender{éis} \\
ellos/ellas & le\ender{en} & comprend\ender{en} \\
\end{tabular}}

\vspace{4pt}

{\footnotesize\textbf{Stamm:} le- / comprend- (Infinitiv ohne -er)}

\end{merkbox}

\vspace{4pt}

\abmark{1--2}

\vspace{8pt}

% ═══════════════════════════════════════════════════════════════
% ABSCHNITT 3: Konjugation -IR
% ═══════════════════════════════════════════════════════════════

\begin{merkbox}[Verbos en -IR: abrir, escribir]

\vspace{2pt}

{\small
\renewcommand{\arraystretch}{1.4}
\begin{tabular}{@{}l|l|l@{}}
& \textbf{abrir} \de{öffnen} & \textbf{escribir} \de{schreiben} \\
\hline
yo & abr\endir{o} & escrib\endir{o} \\
tú & abr\endir{es} & escrib\endir{es} \\
él/ella & abr\endir{e} & escrib\endir{e} \\
nosotros & abr\endir{imos} & escrib\endir{imos} \\
vosotros & abr\endir{ís} & escrib\endir{ís} \\
ellos/ellas & abr\endir{en} & escrib\endir{en} \\
\end{tabular}}

\vspace{4pt}

{\footnotesize\textbf{Unterschied zu -ER:} nosotros \endir{-imos} / vosotros \endir{-ís}}

\end{merkbox}

\vspace{4pt}

\abmark{3--4}

\columnbreak

% ═══════════════════════════════════════════════════════════════
% ABSCHNITT 4: Vergleich der Endungen
% ═══════════════════════════════════════════════════════════════

\begin{grammatikbox}[Vergleich: -AR / -ER / -IR]

\vspace{2pt}

{\small
\renewcommand{\arraystretch}{1.4}
\begin{tabular}{@{}l|ccc@{}}
& \endar{-AR} & \ender{-ER} & \endir{-IR} \\
\hline
yo & \endar{-o} & \ender{-o} & \endir{-o} \\
tú & \endar{-as} & \ender{-es} & \endir{-es} \\
él/ella & \endar{-a} & \ender{-e} & \endir{-e} \\
nosotros & \endar{-amos} & \ender{-emos} & \endir{-imos} \\
vosotros & \endar{-áis} & \ender{-éis} & \endir{-ís} \\
ellos & \endar{-an} & \ender{-en} & \endir{-en} \\
\end{tabular}}

\vspace{6pt}

{\footnotesize
\textbf{Merke:}
\begin{itemize}[leftmargin=1em, itemsep=1pt]
    \item \textbf{yo}: immer \textbf{-o} (alle Konjugationen!)
    \item \ender{-ER} und \endir{-IR}: fast identisch
    \item Unterschied nur bei \textbf{nosotros} und \textbf{vosotros}
\end{itemize}}

\end{grammatikbox}

\vspace{8pt}

% ═══════════════════════════════════════════════════════════════
% ABSCHNITT 5: Vokabeln Schulkontext
% ═══════════════════════════════════════════════════════════════

\begin{vokabelbox}[Verbos en el instituto]

\vspace{2pt}

{\small
\renewcommand{\arraystretch}{1.3}
\begin{tabular}{@{}ll@{}}
\multicolumn{2}{@{}l}{\textbf{Verben auf \endar{-AR}:}} \\[2pt]
\es{hablar} & \de{sprechen} \\
\es{escuchar} & \de{hören, zuhören} \\
\es{mirar} & \de{schauen, ansehen} \\
\es{explicar} & \de{erklären} \\
\es{practicar} & \de{üben, praktizieren} \\[4pt]
\multicolumn{2}{@{}l}{\textbf{Verben auf \ender{-ER}:}} \\[2pt]
\es{leer} & \de{lesen} \\
\es{comprender} & \de{verstehen} \\[4pt]
\multicolumn{2}{@{}l}{\textbf{Verben auf \endir{-IR}:}} \\[2pt]
\es{escribir} & \de{schreiben} \\
\es{abrir} & \de{öffnen} \\
\end{tabular}}

\end{vokabelbox}

\vspace{4pt}

\abmark{5--6}

\vspace{8pt}

% ═══════════════════════════════════════════════════════════════
% ABSCHNITT 6: Beispielsätze
% ═══════════════════════════════════════════════════════════════

\begin{merkbox}[Ejemplos – En clase]

\vspace{2pt}

{\small
\begin{itemize}[leftmargin=1em, itemsep=3pt]
    \item \es{El profesor explica la gramática.}\\
          \de{Der Lehrer erklärt die Grammatik.}
    \item \es{Los alumnos escriben en el cuaderno.}\\
          \de{Die Schüler schreiben ins Heft.}
    \item \es{María lee un libro en español.}\\
          \de{María liest ein Buch auf Spanisch.}
    \item \es{¿Comprendes la tarea?}\\
          \de{Verstehst du die Aufgabe?}
    \item \es{Abrimos el libro en la página 42.}\\
          \de{Wir öffnen das Buch auf Seite 42.}
    \item \es{Practicamos los verbos.}\\
          \de{Wir üben die Verben.}
\end{itemize}}

\end{merkbox}

\end{multicols}

% ═══════════════════════════════════════════════════════════════
% FOOTER
% ═══════════════════════════════════════════════════════════════

\vfill

\noindent\rule{\textwidth}{0.5pt}

\vspace{2pt}

{\footnotesize\textcolor{mittelgrau}{Qué pasa 1 | Verbos regulares: 2ª y 3ª conjugación (-ER / -IR)}}

\end{document}
